\documentclass[11pt,a4paper]{scrbook}

\usepackage[english]{babel}
\usepackage{fontspec}
\usepackage{microtype}
\usepackage[a4paper,left=23mm,right=53mm,top=25mm,bottom=28mm,marginparwidth=38mm,marginparsep=7mm]{geometry}
\usepackage{marginnote}
\usepackage{xcolor}
\usepackage[hidelinks]{hyperref}

\IfFontExistsTF{Libertinus Serif}{%
  \setmainfont{Libertinus Serif}
  \setsansfont{Libertinus Sans}[Scale=0.94]
}{%
  \setmainfont{Latin Modern Roman}
  \setsansfont{Latin Modern Sans}
}
\definecolor{Accent}{HTML}{315B6B}

\setkomafont{chapter}{\sffamily\bfseries\Huge}
\setkomafont{section}{\sffamily\bfseries\Large}

\newcommand{\margininfo}[1]{\marginnote{\footnotesize\sffamily #1}}
\newenvironment{languagenote}
{\par\medskip\noindent{\sffamily\bfseries\color{Accent}Language Note}\par\smallskip}
{\par\medskip}

\begin{document}

\frontmatter
\title{Learner's Guide}
\subtitle{Book/Guide + Kao-inspired information architecture}
\author{Author}
\maketitle
\tableofcontents

\mainmatter

\chapter{Core concept}
\margininfo{Use the margin for supporting information that benefits from proximity.}

The main narrative carries explanation and progression.

\begin{languagenote}
Recurring semantic elements should keep the same typography and spacing throughout the guide.
\end{languagenote}

\section{Example}
A worked example belongs close to the explanation it demonstrates.

\section{Practice}
For instructional books, transition from explanation to example to meaningful practice
without forcing every chapter into an identical sequence when the content does not require it.

\end{document}
